\section{Architecture}
\label{sec:architecture}

\subsection{Four processes, one graph}

The loop decomposes into six responsibilities, and the partitioning rule is
that a component owns exactly one decision, so that changing how that decision
is made touches one place.

\begin{enumerate}[label=\textbf{S\arabic*.},leftmargin=2.8em,itemsep=2pt]
  \item \textbf{Knowledge acquisition.} Turn whatever documents exist into
    queryable, individually identified requirements with provenance back to the
    text they came from. Owner: the ingestion pipeline.
  \item \textbf{Test selection.} Decide what is worth testing next given
    everything currently known. Owner: the planner. This is where the system's
    quality is bounded.
  \item \textbf{Execution.} Drive a running application toward a stated
    objective and record what happened. Owner: the player, one per platform.
  \item \textbf{Evidence reduction.} Turn a trajectory into durable,
    deduplicated claims and decide what the run proved. Owner: the
    investigator.
  \item \textbf{Persistence and retrieval.} Hold the knowledge and answer
    questions about it. Owner: the retrieval service, which is the only writer
    to the database.
  \item \textbf{Operation.} Let a person start a campaign, understand a verdict
    and take the result away. Owner: the gateway and the dashboard.
\end{enumerate}

These map onto four concurrent processes: the agent gateway on port 9100, which
is the only entry point and holds the planner; the retrieval service on port
9010, which is the only process that opens a connection to the graph database;
the player, which is a campaign driver rather than a server; and Neo4j itself
on port 7687.

\begin{sidewaysfigure}
  \centering
  \includegraphics[width=0.96\textheight]{01-system-architecture.png}
  \caption[System architecture]{System architecture. Amber is an input, green
  is knowledge, blue is a service, purple is one of the three agent roles, red
  is a language-model call, grey is the device. Every input is optional: the
  system degrades to whatever sources exist rather than failing without them.}
  \label{fig:architecture}
\end{sidewaysfigure}

\subsection{Three decisions that fix the shape}

\subsubsection{The graph is the only channel between roles}

The planner, the player and the investigator never call one another. Each is a
stateless language-model call that reads from and writes to the same graph, so
the system's memory sits in one inspectable place and no conversational state
is carried between rounds.

This is what makes learning a property that can be audited rather than an
emergent property of a long conversation. After a round, the graph either grew
or it did not, and the difference is a query. It also means any role can be
replaced without touching the others, which is what made the platform work of
\Cref{sec:capabilities} a matter of writing a second player.

\subsubsection{One writer}

Only the retrieval service opens a Neo4j connection. The gateway and the player
reach the graph exclusively through its HTTP interface. The alternative, giving
each process its own driver, would have spread the schema across three
codebases and turned a constraint change into a three-way migration. It also
gives one place to enforce the write invariants that keep the application model
consistent under concurrent runs.

\subsubsection{The planner never writes steps}

The planner's output is an objective in plain language, not a tap sequence,
because it cannot see the live screen. A wrong tap script is worse than an
imprecise objective: the player can recover from a vague goal by looking at
what is actually in front of it, and cannot recover from instructions that do
not match the screen.

\subsection{The knowledge model}

Two graphs coexist in one store. An \emph{a priori} graph holds what the
application should do, built from documents: requirements, specification
chunks, validation rules. An \emph{a posteriori} graph holds what it actually
does, built from observation: interface states, transitions, navigation paths,
execution logs. Test cases and findings span the two, and that span is the
central idea of the design. Exploration is steered by the gap between what the
specification promises and what the running application was observed to do.

\begin{sidewaysfigure}
  \centering
  \includegraphics[width=0.96\textheight]{03-knowledge-graph-schema.png}
  \caption[Knowledge graph schema]{The knowledge graph. Specification-derived
  nodes and observation-derived nodes share one store, joined by the test cases
  and findings that span them.}
  \label{fig:schema}
\end{sidewaysfigure}

\subsection{Ingestion: documents to queryable requirements}

Ingestion decides what the rest of the system can ever know, which is why it is
treated as a first-class component rather than as a loader. A document is
chunked, embedded, and reduced to individually identified requirements and
validation rules, each linked back to the chunk it came from. Any later claim
can therefore be traced to the sentence that justified it, which is what makes
requirement coverage reportable at all rather than merely assertable.

Re-ingesting a changed specification produces a new version and a drift report
rather than overwriting the previous one, so a requirement that disappears from
a document does not silently disappear from the coverage denominator.

\begin{figure}[H]
  \centering
  \includegraphics[width=\textwidth]{11-ingestion-document-to-queryable-requirements.png}
  \caption[Ingestion pipeline]{Prose becomes atomic, identified requirements
  with provenance back to their source chunk.}
  \label{fig:ingestion}
\end{figure}

Ingestion is format-agnostic. PDF, Office documents, HTML and Markdown are all
accepted, and the pipeline auto-detects better parsers when they are installed
without requiring them. Embeddings run locally through ONNX, so retrieval
carries no per-call charge and no dependency on an external service.

\subsection{One round, end to end}

\begin{figure}[H]
  \centering
  \includegraphics[width=0.97\textwidth]{02-one-test-round-end-to-end.png}
  \caption[One test round]{One test round. The planner reads the graph,
  proposes, and is validated before anything reaches a device. The player
  drives up to its interaction budget. The investigator reduces the trajectory
  to findings, and is given the findings already known for those screens so
  that a repeat observation reinforces an existing finding instead of creating
  a duplicate.}
  \label{fig:round}
\end{figure}

\subsection{The planner}

The planner answers one question: given everything the graph currently holds,
what is worth testing next? It assembles its own prompt from retrieved
knowledge, proposes a test case with a stated objective and the requirements it
claims to cover, and passes through a validation gate before anything reaches a
device.

\begin{figure}[H]
  \centering
  \includegraphics[width=0.97\textwidth]{05-planner-pipeline-mode-default.png}
  \caption[Planner pipeline]{The default planner as an explicit state machine:
  bootstrap context, plan retrieval, execute retrieval, generate, check for
  duplicates. Each node is separately testable and the flow is inspectable.}
  \label{fig:planner}
\end{figure}

Retrieval draws on six sources: the specification, a design export of screens,
a design export of flows, the observed live interface, defect history, and the
navigation tree. Each is optional. A source that is absent contributes nothing
rather than blocking the round, which is the mechanism behind DR-01.

Two planner modes exist behind one output contract. The default assembles a
fixed context and issues retrieval calls in up to three rounds. The alternative
gives the planner a set of graph tools and lets it investigate the graph
directly, deciding for itself what to look at next.

\begin{figure}[H]
  \centering
  \includegraphics[width=0.97\textwidth]{13-why-tools-replaced-retrieval-sources.png}
  \caption[Why a tool interface was added]{The structural difference between a
  fixed context block and a tool interface. A fixed block must guess in advance
  what the planner will need.}
  \label{fig:tools-why}
\end{figure}

\begin{figure}[H]
  \centering
  \includegraphics[width=0.97\textwidth]{14-the-nine-tools-as-an-investigation.png}
  \caption[The planner's graph tools]{The tool-calling planner treats test
  selection as an investigation over the graph rather than as a single
  generation from a pre-assembled context.}
  \label{fig:tools}
\end{figure}

The tool-calling planner is better formed on every property the team values,
and it is nevertheless not the default. No campaign had been run that showed it
produced better campaigns, and being better designed is not the same as
producing better results. It remains in the source tree behind a configuration
switch, and the comparison is named in \Cref{sec:next} as work still to do.

\subsection{The player}

The player is the only role that sees the application. It receives an objective
in plain language and drives the interface toward it, one observation and one
action at a time, until the objective is met, the objective is refuted, or a
bound is reached.

\begin{figure}[H]
  \centering
  \includegraphics[width=0.97\textwidth]{16-the-executor-in-detail.png}
  \caption[The player in detail]{How a plain-language objective becomes a
  sequence of grounded device actions, and how each step is recorded.}
  \label{fig:executor}
\end{figure}

On Android the player drives the device through the accessibility service using
a companion on-device application that also supplies an input method, which is
what makes text entry into multi-field forms reliable. It reads an
accessibility tree averaging 4{,}686 bytes per step together with a
$1080\times1920$ screenshot, and grounds its action in both.

On the web the same role runs against an accessibility-annotated browser
snapshot rather than raw markup, so controls are addressed by role and name on
both platforms and the planner does not need to know which platform it is
driving. The web player needs no vision model, because the snapshot already
carries the structure a screenshot would have to be read for.

Controls are addressed by accessible role, name and content description rather
than by pixel coordinates. This constrains the player to what the accessibility
layer exposes, which is a real limitation and also the reason its
element-location failures are meaningful as an accessibility signal about the
application under test.

\subsection{The investigator}

The investigator reads a completed trajectory against what the test intended
and reduces it to durable findings. It runs off the critical path, so latency
can be traded for a second opinion, and it is deliberately configured to a
different model family from the players so that no model grades its own
execution.

\begin{figure}[H]
  \centering
  \includegraphics[width=0.97\textwidth]{06-investigator-trajectory-to-knowledge.png}
  \caption[Trajectory to knowledge]{How a raw trajectory becomes deduplicated,
  attributed findings that the next planning round can read.}
  \label{fig:investigator}
\end{figure}

It is given the findings already recorded for the screens the run visited, so a
repeated observation reinforces an existing finding rather than creating a
near-duplicate. An open question that has been probed three times without
resolution is closed as inconclusive, which is a result a human can act on
rather than an endless retry.

\subsection{State identity}

The most consequential calculation in the system is not a budget but a
similarity decision: is this screen one the application model already contains?
Merging too eagerly collapses the map. Splitting too eagerly destroys it. One
wizard, whose form fields carried their current values in the same attribute
that carries a button's label, produced nine distinct states before the cascade
below was designed.

\begin{figure}[H]
  \centering
  \includegraphics[height=0.78\textheight,keepaspectratio]{17-state-identity-is-this-a-screen-we-have-seen.png}
  \caption[State identity cascade]{The five-stage cascade, cheapest test first.
  Each stage was added in response to a specific way the previous one failed,
  and every threshold is configurable.}
  \label{fig:state-identity}
\end{figure}

The stages are ordered by cost. An exact signature match is a string
comparison. Structural Jaccard tolerates dynamic chrome such as counters and
timestamps. Containment recognises a screen captured before its late-painting
chrome arrived, which is a strict subset of the same screen fully rendered and
which Jaccard therefore scores as different; it is guarded by a minimum control
count and a maximum size ratio, because a one-control blank screen is contained
in everything. Skeleton containment compares only resource identifier and
class, asking whether the same widgets are present regardless of their values,
with a floor on full containment so that two different screens built from
similar widget types stay apart. A perceptual hash is the last resort and the
only signal available when the accessibility tree is too thin to describe the
screen at all.

\keypoint{This cascade is what allows the system to work against a target that
exposes no stable identifiers. It is also the component with the most
significant outstanding limitation, reported in \Cref{sec:limits}.}

\subsection{Fault attribution}

A run that does not pass is ambiguous between three explanations, and
conflating them inflates the reported defect count, which is the most expensive
error this system can make.

\begin{figure}[H]
  \centering
  \includegraphics[width=\textwidth]{09-failure-attribution.png}
  \caption[Fault attribution]{Every non-passing execution is classified at the
  point of failure as application evidence, an agent limitation, or an
  environment fault.}
  \label{fig:attribution}
\end{figure}

The classification has exactly one definition in the source tree, so an agent
limitation cannot drift into being reported as an application defect through a
second implementation. The exported run report filters agent-caused failures
before a defect is ever named, so what reaches a stakeholder is the
application's behaviour rather than the tester's.

A degradation register records every fallback and every silent substitution
across process boundaries, which is what allows a campaign to claim zero silent
fallbacks and have the claim be checkable rather than asserted.

\subsection{Explicit bounds}

An exploratory agent has no natural stopping point and will consume whatever it
is given, so every bound is explicit and read from configuration.

\begin{table}[H]
  \centering\small
  \setlength{\tabcolsep}{5pt}
  \begin{tabularx}{\textwidth}{@{}>{\hsize=0.70\hsize}Y>{\hsize=0.42\hsize}Y>{\hsize=1.88\hsize}Y@{}}
    \toprule
    Bound & Value & Why it exists \\
    \midrule
    Device actions per test & 50 & The player's interaction budget. Exhausting it yields no evidence about the application, so the planner is told the number and warned when recent tests repeatedly hit it \\
    \addlinespace
    Wall clock per test & 900\,s & Bounds one execution including model latency and recovery. This is the campaign value; the shipped default is 420\,s \\
    \addlinespace
    Planner turns & 12 & Ceiling on the tool loop \\
    \addlinespace
    Forced proposal after & 6 turns & A model left free will investigate to the ceiling and never commit. Observed on the first live run: ten turns of tool calls and no proposal \\
    \addlinespace
    Proposal rejections & 3 & After three the best effort is accepted and a degradation recorded, so a planning round never returns nothing \\
    \addlinespace
    Tool result caps & 3000 / 2500 / 1200 / 2000 characters & Requirements, findings, one screen, anything else. Each tool is responsible for not flooding the context on its own \\
    \addlinespace
    Finding attempts & 3 & After three probes an open question auto-closes as inconclusive \\
    \addlinespace
    Knowledge half-life & 90 days & Decays the weight of stale observations so the agent adapts when the application changes \\
    \bottomrule
  \end{tabularx}
\end{table}
